\chapter{Format of This Cheatsheet}

I personally believe having a consistent format would help aid readers
in obtaining information regarding a certain topic, and will also help the author(me)
in writing concepts without missing certain details within a concept:

\vspace{3em}

Each new introduction of a \textbf{new concept} will have the following:
\begin{enumerate}
    \item A hopefully \textbf{detailed overview} of the concept in general.
    \item The \textbf{technical definition / mathematical term} of such a concept.
    \item \textbf{Example use-cases and explanation} of the use-case of such a concept.
\end{enumerate}

\tgap

\emph{And \textbf{optionally}}:
\begin{enumerate}
    \item Some relevant / interesting \textbf{notes} regarding the concept.
    \item Additional resources as \textbf{footnotes} to explain the concept in more detail.
    \item If the mathematical definition has fancy letters and symbols, it will get explained
        in the "\textbf{where}" section right after the definition.
    \item \textbf{Counter-examples and explaination} of why it doesn't work this way.
        (It's optional because sometimes counter-examples are harder to explain)
\end{enumerate}

\tgap

\emph{For \textbf{text formatting}}:
\begin{enumerate}
    \item Important / Relevant new words for concepts will get their fonts \textbf{bolded}.
\end{enumerate}



